\section{Averaging as an Approximate Weighted LMO}
\label{app:matnormal-p2}

The single-batch LMO in \Cref{sec:sample-loss-lmo} has two inputs, the
gradient in its objective and the per-sample second moment in its
outergradient constraint. To draw on recent batches while keeping the LMO
local to the current iterate, we evaluate every recent sample's gradient at
$W_t$ and average the two inputs across those samples, giving the ideal
moments $\overline G_t$ and $\overline\Sigma_t$. This appendix explains how
\Cref{alg:ours}
approximates them from the LoRA factor gradients already computed by
backpropagation.

\Cref{sec:weighted-lmo-p2} writes the weighted LMO and isolates the two
moments through which past steps enter the problem.
\Cref{sec:weighted-estimation-p2}
replaces those moments by quantities built from $G_{A,s}$ and $G_{B,s}$,
giving the $p,q$ recursion in \Cref{alg:ours} and the
update~\eqref{eq:polora-update}.
\Cref{sec:scale-placement-p2} explains why \Cref{alg:ours} normalizes the
accumulated vectors after averaging rather than normalizing each estimate
before it enters the average. The update is invariant to the scale of the
curvature factors, but the average that estimates them depends on that scale.

\subsection{The Weighted LMO}
\label{sec:weighted-lmo-p2}

Let $s$ index optimizer steps, and let $i$ index samples in the
minibatch of size $n_s$ at step $s$. Write
$G_{s,i}:=\nabla\ell_{x_{s,i}}(W_s)$,
$g_{s,i}:=\vec(G_{s,i})$, and
$G_s:=n_s^{-1}\sum_iG_{s,i}$. Let $m_{t,s}$ and $w_{t,s}$ be
nonnegative weights on past steps, one for the objective and one for
the constraints.

In the ideal weighted LMO, every sample from a step $s\le t$ contributes
the gradient it would have at the current weights $W_t$, so drawing on past
steps changes only which samples enter the averages, not where their
gradients are evaluated. The two resulting moments are
\[
\overline G_t
  =
  \sum_{s\le t}\frac{m_{t,s}}{n_s}\sum_{i=1}^{n_s}
  \nabla\ell_{x_{s,i}}(W_t),
\qquad
\overline\Sigma_t
  =
  \sum_{s\le t}\frac{w_{t,s}}{n_s}\sum_{i=1}^{n_s}
  \vec(\nabla\ell_{x_{s,i}}(W_t))
  \vec(\nabla\ell_{x_{s,i}}(W_t))^\top .
\]
The first moment is the gradient in the linear objective. For the
constraint, the relevant gradients are rescaled as
$\sqrt{w_{t,s}/n_s}\,\nabla\ell_{x_{s,i}}(W_t)$, and their outer
products sum to $\overline\Sigma_t$. Applying the leverage argument of
\Cref{lem:gradprop} to these rescaled gradients gives a relaxation of
the outergradient constraint with the same form
as~\eqref{eq:polorafull-hard}. The relaxed problem is
\begin{equation}\label{eq:weighted-relaxed-p2}
 \begin{array}{ll}
  \underset{\Delta W}{\mathrm{minimize}} &
  \ip{\overline G_t, \Delta W} \\[2pt]
  \mbox{subject to} &
  \displaystyle
  \max_{\widetilde{G}\in\mathcal{G}(\overline\Sigma_t)}
  \ip{\widetilde{G},\Delta W}\le\tau .
  \end{array}
\end{equation}
\Cref{alg:ours} stores neither directly, because it never recomputes old
samples' gradients at $W_t$ and the per-sample second moment is too large
to keep.

\paragraph{The exponential weights.}
If every past gradient could be recomputed at $W_t$, equal weights would be
the direct extension of the single-batch LMO to samples from earlier steps.
Since that average is unavailable, \Cref{alg:ours} uses two exponentially
decaying weight sequences, one for the objective and one for the constraint.
These weights use constant memory and put less mass on older gradients and
factors, which were computed farther from the current iterate.

\Cref{alg:ours} uses look-ahead momentum for the first moment,
\[
  M_{t+1}=\beta_1M_t+(1-\beta_1)G_t,
  \qquad
  \widehat M_t=\beta_1M_{t+1}+(1-\beta_1)G_t,
\]
and an EMA for the second moment. Unrolling these recursions
gives the induced step weights
\begin{equation}\label{eq:EMAweights-p2}
  m_{t,s}=\beta_1^2(1-\beta_1)\beta_1^{t-1-s}
  \quad(s<t),
  \qquad
  m_{t,t}=1-\beta_1^2,
  \qquad
  w_{t,s}=(1-\beta_2)\beta_2^{t-s}.
\end{equation}
We write $\EMA_{s\le t}[\phi(G_s)]=\sum_{s\le t}w_{t,s}\,\phi(G_s)$ for an
average over past steps with weights $w_{t,s}$. Replacing each unavailable
current-weight gradient in $\overline G_t$ by the observed batch gradient
$G_s$ turns the first moment into the look-ahead momentum $\widehat M_t$,
with the objective weights $m_{t,s}$ in~\eqref{eq:EMAweights-p2}. The
optimizer never forms the $d_{\mathrm{out}}\times d_{\mathrm{in}}$ matrix
$\widehat M_t$, so the next subsection expresses its needed projections and
the curvature EMA in terms of factor gradients.

\subsection{Estimating the Moments from Factor Gradients}
\label{sec:weighted-estimation-p2}

The optimizer maintains momenta for $G_{A,s},G_{B,s}$ and diagonal
curvature vectors instead of the dense moments $\overline G_t$ and
$\overline\Sigma_t$.

\paragraph{From moments at $W_t$ to momentum and curvature variables.}
The ideal LMO in~\eqref{eq:weighted-relaxed-p2} is written with old samples
re-evaluated at $W_t$. \Cref{alg:ours} replaces those ideal moments by stored
variables through the following modeling steps. Only the first step is a
stale-gradient approximation; the remaining steps concern the observed recent
gradient stream and the changing LoRA factor maps.
\begin{enumerate}[leftmargin=1.6em,itemsep=2pt,topsep=2pt]
  \item A gradient of an old sample at $W_t$ is replaced by the gradient
  observed when that sample was processed,
  $\nabla\ell_{x_{s,i}}(W_t)\leadsto\nabla\ell_{x_{s,i}}(W_s)$.
  The replacement uses separate weighted stale-gradient assumptions for the
  objective and the curvature. For the objective weights $m_{t,s}$, assume
  \[
    \left\|
      \sum_{s\le t}\frac{m_{t,s}}{n_s}\sum_{i=1}^{n_s}
      \bigl(\nabla\ell_{x_{s,i}}(W_t)-\nabla\ell_{x_{s,i}}(W_s)\bigr)
    \right\|_2
    =
    O(\epsilon)
    \left\|
      \sum_{s\le t}\frac{m_{t,s}}{n_s}\sum_{i=1}^{n_s}
      \nabla\ell_{x_{s,i}}(W_s)
    \right\|_2 .
  \]
  For the curvature weights $w_{t,s}$, write
  $h^{(t)}_{s,i}=\vec(\nabla\ell_{x_{s,i}}(W_t))$ and
  $h^{(s)}_{s,i}=\vec(\nabla\ell_{x_{s,i}}(W_s))$, and assume
  \[
    \left\|
      \sum_{s\le t}\frac{w_{t,s}}{n_s}\sum_{i=1}^{n_s}
      \bigl(h^{(t)}_{s,i}(h^{(t)}_{s,i})^\top
      -h^{(s)}_{s,i}(h^{(s)}_{s,i})^\top\bigr)
    \right\|_2
    =
    O(\epsilon)
    \left\|
      \sum_{s\le t}\frac{w_{t,s}}{n_s}\sum_{i=1}^{n_s}
      h^{(s)}_{s,i}(h^{(s)}_{s,i})^\top
    \right\|_2 .
  \]
  The rest of the section works with the observed gradients $G_s$.
  \item The second moment within the batch is approximated by the outer
  product of the batch gradient,
  \[
    \frac1{n_s}\sum_i g_{s,i}g_{s,i}^\top
    \quad\leadsto\quad
    g_sg_s^\top,
    \qquad g_s=\vec(G_s).
  \]
  For $n_s>1$, this replaces the per-sample second moment by the
  second moment of the minibatch gradient, as in AdaGrad~\cite{adagrad},
  Adam~\cite{adam}, and Shampoo~\cite{shampoo}.
  \item When relating the dense observed moment to factor-gradient moments,
  replacing $B_s^\top G_s$ by $B_t^\top G_s$ and
  $G_sA_s^\top$ by $G_sA_t^\top$ assumes weighted factor-projection
  stability. The concrete stability bounds used by the online fit are stated in
  \Cref{prop:online-stationary-p2}; it measures the whole quadratic-form
  change, including the cross terms created by changing the factors.
  \item The resulting local second moment is modeled exactly by diagonal
  Kronecker factors $P^\star_t,Q^\star_t$,
  \begin{equation}\label{e-diag-kron-second-mom-model-p2}
    \EMA_{s\le t}[g_sg_s^\top]
    =
    \kappa_t(Q^\star_t\otimes P^\star_t).
  \end{equation}
  We fix $\norm{\diag P^\star_t}_\infty=\norm{\diag Q^\star_t}_\infty=1$ and
  let the scalar $\kappa_t>0$ carry the magnitude removed by the final
  normalization.
\end{enumerate}
The stale-gradient and batch-gradient approximations reduce the constraint
moment to
\[
  \overline\Sigma_t\approx\EMA_{s\le t}[g_sg_s^\top],
\]
which the diagonal Kronecker model~\eqref{e-diag-kron-second-mom-model-p2}
writes as $\kappa_t(Q^\star_t\otimes P^\star_t)$.

\paragraph{Factor momenta.}
LoRA backpropagation already computes the factor gradients
$G_{A,s}=B_s^\top G_s$ and $G_{B,s}=G_sA_s^\top$. Forming the dense
gradient $G_s$ would require a $d_{\mathrm{out}}\times d_{\mathrm{in}}$
matrix per layer. The optimizer therefore represents the first moment only
through the two projections needed by the factor update. Under the
linearized product update~\eqref{eq:DeltaW-lin}, the objective of
\eqref{eq:weighted-relaxed-p2} decomposes as
\[
  \ip{\widehat M_t,\,\Delta B\,A+B\,\Delta A}
  =
  \ip{\widehat M_tA^\top,\Delta B}
  +
  \ip{B^\top\widehat M_t,\Delta A}.
\]
The direction step then needs only $\widehat M_tA_t^\top$ and
$B_t^\top\widehat M_t$. The factor momenta in \Cref{alg:ours} average
$G_sA_s^\top$ and $B_s^\top G_s$ with the objective weights $m_{t,s}$.
Replacing $A_s,B_s$ by $A_t,B_t$ in these first-moment sums requires the
corresponding $m$-weighted projection-replacement bounds
\[
\begin{aligned}
  \left\|
    \sum_{s\le t}m_{t,s}(B_s^\top G_s-B_t^\top G_s)
  \right\|_2
  &=
  O(\epsilon)
  \left\|
    \sum_{s\le t}m_{t,s}B_t^\top G_s
  \right\|_2,\\
  \left\|
    \sum_{s\le t}m_{t,s}(G_sA_s^\top-G_sA_t^\top)
  \right\|_2
  &=
  O(\epsilon)
  \left\|
    \sum_{s\le t}m_{t,s}G_sA_t^\top
  \right\|_2 .
\end{aligned}
\]
The $w$-weighted bounds above are used for the curvature estimator.

\paragraph{Moments of the factor gradients.}
With $\EMA_{s\le t}[g_sg_s^\top]$ as the dense second moment, replacing
$A_s,B_s$ by $A_t,B_t$ links the second moments of the factor gradients to
the weighted second moment of $G_s$. For the $A$ factor,
$\vec(G_{A,s})=(I\otimes B_s^\top)\vec(G_s)$, so
\begin{align*}
 \Sigma_{A,t}
 &:= \EMA_{s\le t}\big[\vec(G_{A,s})\vec(G_{A,s})^\top\big] \\
 &= \EMA_{s\le t}\big[
      (I\otimes B_s^\top)g_sg_s^\top(I\otimes B_s)
    \big] \\
 &\approx
    (I\otimes B_t^\top)
    \EMA_{s\le t}\big[g_sg_s^\top\big]
    (I\otimes B_t) \\
 &\approx
    \kappa_t\,(I\otimes B_t^\top)
    (Q^\star_t\otimes P^\star_t)
    (I\otimes B_t) \\
 &= \kappa_t\,Q^\star_t\otimes(B_t^\top P^\star_tB_t).
\end{align*}
The first approximation replaces $B_s$ by $B_t$ inside the weighted
sum, the second applies the diagonal Kronecker model, and the final
equality is the mixed product rule. The analogous calculation for
$G_{B,s}$, written with $\vec(G_{B,s}^\top)$ so that the Kronecker
blocks are indexed by the output coordinate, gives
\begin{align}
  \Sigma_{A,t}
    &\approx \kappa_t\,Q^\star_t\otimes C^\star_{A,t},
  \nonumber \\
  \Sigma_{B,t}
    &:= \EMA_{s\le t}\big[
      \vec(G_{B,s}^\top)\vec(G_{B,s}^\top)^\top
    \big]
    \approx \kappa_t\,P^\star_t\otimes C^\star_{B,t},
  \label{eq:observable-moments-p2}
\end{align}
where
$C^\star_{A,t}=B_t^\top P^\star_tB_t$ and
$C^\star_{B,t}=A_tQ^\star_tA_t^\top$.

\paragraph{Fitting the diagonal factors.}
We fit the diagonal factors to the factor-gradient moments in
\eqref{eq:observable-moments-p2}. In this paragraph, all quantities are
at the current step, so we suppress the subscript $t$. Following
KL-Shampoo~\cite{klshampoo}, the fit uses the log-determinant
divergence~\cite{logdetdiv}
\[
  D(\Sigma,S)=\Tr(\Sigma S^{-1})-\log\det(\Sigma S^{-1})-d,
\]
defined for positive definite matrices of common dimension $d$.
Under the block model below, the $A$-gradient moment would recover
$Q^\star$ once a small $r\times r$ closure matrix is fixed, and the
$B$-gradient moment would recover $P^\star$ once its small closure matrix is
fixed. The fit therefore freezes the opposite-side closure:
\begin{equation}\label{eq:factor-fit-p2}
  \hat{q}(\widehat C_A)
  =
  \argmin_{q>0}
  D\bigl(\Sigma_A,\diag(q)\otimes\widehat C_A\bigr),
  \qquad
  \hat{p}(\widehat C_B)
  =
  \argmin_{p>0}
  D\bigl(\Sigma_B,\diag(p)\otimes\widehat C_B\bigr),
\end{equation}
where $\widehat C_A,\widehat C_B\succ0$. The blockwise first-order
conditions give the closed forms
\[
  \hat{q}(\widehat C_A)
  =
  \frac1r\diag\!\left(
    \EMA_{s\le t}[G_{A,s}^\top\widehat C_A^{-1}G_{A,s}]
  \right),
  \qquad
  \hat{p}(\widehat C_B)
  =
  \frac1r\diag\!\left(
    \EMA_{s\le t}[G_{B,s}\widehat C_B^{-1}G_{B,s}^\top]
  \right).
\]

For fixed positive definite $\Sigma$ and $S$, define the best scalar
multiple
\[
  c^\star(\Sigma,S)
  :=
  \argmin_{c>0}D(\Sigma,cS)
  =
  \Tr(\Sigma S^{-1})/d .
\]

\begin{lemma}[Closure invariance of the fixed fit]\label{lem:consistency-p2}
Assume the moments of the factor gradients satisfy
\[
  \Sigma_A=\kappa\,Q^\star\otimes C^\star_A,
  \qquad
  \Sigma_B=\kappa\,P^\star\otimes C^\star_B
\]
exactly, with
$C^\star_A=B^\top P^\star B$,
$C^\star_B=AQ^\star A^\top$, and
$\norm{\diag P^\star}_\infty=\norm{\diag Q^\star}_\infty=1$, with all
diagonal entries of $P^\star,Q^\star$ positive. Assume also that
$C^\star_A$ and $C^\star_B$ are positive definite.
For any positive definite closures $\widehat C_A$ and $\widehat C_B$, the
fits~\eqref{eq:factor-fit-p2} return
\[
  \hat{q}(\widehat C_A)
  =
  \kappa\,c^\star(C^\star_A,\widehat C_A)\,\diag(Q^\star),
  \qquad
  \hat{p}(\widehat C_B)
  =
  \kappa\,c^\star(C^\star_B,\widehat C_B)\,\diag(P^\star).
\]
Therefore
$\hat{q}/\norm{\hat{q}}_\infty=\diag(Q^\star)$ and
$\hat{p}/\norm{\hat{p}}_\infty=\diag(P^\star)$.
\end{lemma}
\begin{proof}
Under the stated moment model, $\Sigma_A$ is block diagonal
with $j$th block $\kappa q^\star_j C^\star_A$, while
$\diag(q)\otimes\widehat C_A$ has $j$th block $q_j\widehat C_A$.
Traces and log-determinants add across blocks, so the objective for
$q$ is
\[
  \sum_j D(\kappa q^\star_j C^\star_A, q_j\widehat C_A).
\]
The $j$th term is minimized by its best scalar multiple,
\[
  \hat q_j
  =
  c^\star(\kappa q^\star_j C^\star_A,\widehat C_A)
  =
  \kappa q^\star_j c^\star(C^\star_A,\widehat C_A).
\]
The scalar $\kappa c^\star(C^\star_A,\widehat C_A)$ is independent of
$j$, and the normalization removes it. The proof for $\hat p$ is the
same with the two factors exchanged.
\end{proof}

\paragraph{Online estimator.}
The weighted fit in~\eqref{eq:factor-fit-p2} would combine all past
factor gradients using the same current closures. The online recurrence
instead evaluates each damped observation once, using the closures available
at that step. Given raw EMA vectors $p_t,q_t$, the implementation forms
the effective diagonal metrics
\[
  D_P(p_t)=\diag\!\bigl(p_t/\norm{p_t}_\infty+\delta\mathbf 1\bigr),
  \qquad
  D_Q(q_t)=\diag\!\bigl(q_t/\norm{q_t}_\infty+\delta\mathbf 1\bigr),
\]
where $\delta>0$ is the relative damping parameter. The raw vectors
$p_t,q_t$ are initialized with positive entries, so these normalizations are
defined before the first curvature update. We state the inverse formulas for
nondegenerate steps where the damped closures below are positive definite:
\[
\begin{aligned}
  C_{A,t}&=B_t^\top D_P(p_t)B_t,
  &
  C_{B,t}&=A_tD_Q(q_t)A_t^\top,\\
  \widetilde C_{A,t}^{-1}
  &=(C_{A,t}+\delta\lambda_{\max}(C_{A,t})I)^{-1},
  &
  \widetilde C_{B,t}^{-1}
  &=(C_{B,t}+\delta\lambda_{\max}(C_{B,t})I)^{-1}.
\end{aligned}
\]
The Newton--Schulz routine computes the inverse square roots whose squares
give these inverse closures, up to numerical approximation. The online
recurrence is
\begin{equation}\label{eq:ours-klcoupling-p2}
\begin{aligned}
  \hat{q}_t
    &= \frac1r\diag\!\left(G_{A,t}^\top \widetilde C_{A,t}^{-1}G_{A,t}\right),
  &
  \hat{p}_t
    &= \frac1r\diag\!\left(G_{B,t}\widetilde C_{B,t}^{-1}G_{B,t}^\top\right),
  \\
  q_{t+1}
    &= \beta_2q_t+(1-\beta_2)\hat{q}_t,
  &
  p_{t+1}
    &= \beta_2p_t+(1-\beta_2)\hat{p}_t,
  \\
  Q_{t+1}
    &= \diag\!\bigl(q_{t+1}/\norm{q_{t+1}}_\infty\bigr),
  &
  P_{t+1}
    &= \diag\!\bigl(p_{t+1}/\norm{p_{t+1}}_\infty\bigr).
\end{aligned}
\end{equation}
Unrolling the EMA gives
\[
  q_{t+1}
  =
  \beta_2^{t+1}q_0+q^\circ_{t+1},
  \qquad
  p_{t+1}
  =
  \beta_2^{t+1}p_0+p^\circ_{t+1},
  \qquad
  q^\circ_{t+1}=\sum_{s\le t}w_{t,s}\hat q_s,
  \qquad
  p^\circ_{t+1}=\sum_{s\le t}w_{t,s}\hat p_s .
\]
The proposition below compares these tail-free averages with the fitted
vectors obtained by freezing the factor maps and inverse closures at the
current model factors. The raw vectors $q_{t+1},p_{t+1}$ differ from
$q^\circ_{t+1},p^\circ_{t+1}$ by the displayed initialization terms.

\begin{proposition}[EMA tracking from weighted factor and closure stability]
\label{prop:online-stationary-p2}
At step $t$, let
$q^\star_t=\diag(Q^\star_t)$ and
$p^\star_t=\diag(P^\star_t)$ be positive vectors with
$\norm{q^\star_t}_\infty=\norm{p^\star_t}_\infty=1$, defined by the exact
dense moment model
\[
  \sum_{s\le t}w_{t,s}\,g_sg_s^\top
  =
  \kappa_t(Q^\star_t\otimes P^\star_t).
\]
Set
$C^\star_{A,t}=B_t^\top P^\star_tB_t$ and
$C^\star_{B,t}=A_tQ^\star_tA_t^\top$, and assume these closures are
positive definite. Define
\[
  \overline G_{A,s}=B_t^\top G_s,\qquad
  \overline G_{B,s}=G_sA_t^\top,\qquad
  H^A_t=(C^\star_{A,t})^{-1},\qquad
  H^B_t=(C^\star_{B,t})^{-1}.
\]
Let $K^A_s=\widetilde C_{A,s}^{-1}$ and
$K^B_s=\widetilde C_{B,s}^{-1}$. Assume the observed factor maps and the
online inverse closures are stable in the weighted quadratic forms
\[
\begin{aligned}
  \frac{1}{r\kappa_t}
  \left\|
    \sum_{s\le t}w_{t,s}
    \bigl(
      G_{A,s}^\top H^A_tG_{A,s}
      -\overline G_{A,s}^\top H^A_t\overline G_{A,s}
    \bigr)
  \right\|_2
  &=
  O(\epsilon),
  \\
  \frac{1}{r\kappa_t}
  \left\|
    \sum_{s\le t}w_{t,s}
    G_{A,s}^\top(K^A_s-H^A_t)G_{A,s}
  \right\|_2
  &=
  O(\epsilon),
  \\
  \frac{1}{r\kappa_t}
  \left\|
    \sum_{s\le t}w_{t,s}
    \bigl(
      G_{B,s}H^B_tG_{B,s}^\top
      -\overline G_{B,s}H^B_t\overline G_{B,s}^\top
    \bigr)
  \right\|_2
  &=
  O(\epsilon),
  \\
  \frac{1}{r\kappa_t}
  \left\|
    \sum_{s\le t}w_{t,s}
    G_{B,s}(K^B_s-H^B_t)G_{B,s}^\top
  \right\|_2
  &=
  O(\epsilon).
\end{aligned}
\]
Then, for $\epsilon$ below a fixed numerical constant,
\[
\begin{aligned}
  \left\|
    \frac{q^\circ_{t+1}}{\norm{q^\circ_{t+1}}_\infty}
    -
    q^\star_t
  \right\|_\infty
  &=
  O(\epsilon),
  \\
  \left\|
    \frac{p^\circ_{t+1}}{\norm{p^\circ_{t+1}}_\infty}
    -
    p^\star_t
  \right\|_\infty
  &=
  O(\epsilon).
\end{aligned}
\]
\end{proposition}
\begin{proof}
Define the ideal frozen fits
\[
\begin{aligned}
  q^{\mathrm{id}}_t
  &=
  \frac1r\diag\!\sum_{s\le t}w_{t,s}
  \overline G_{A,s}^\top H^A_t\overline G_{A,s},
  \\
  p^{\mathrm{id}}_t
  &=
  \frac1r\diag\!\sum_{s\le t}w_{t,s}
  \overline G_{B,s}H^B_t\overline G_{B,s}^\top .
\end{aligned}
\]
The exact dense moment model and the mixed-product rule give
$q^{\mathrm{id}}_t=\kappa_tq^\star_t$ and
$p^{\mathrm{id}}_t=\kappa_tp^\star_t$. Adding and subtracting the terms with
$H^A_t$ and $H^B_t$ in the online averages, the four stability bounds and
$\norm{\diag M}_\infty\le\norm{M}_2$ imply
\[
  \norm{q^\circ_{t+1}-q^{\mathrm{id}}_t}_\infty
  +
  \norm{p^\circ_{t+1}-p^{\mathrm{id}}_t}_\infty
  =
  O(\epsilon)\kappa_t .
\]
For any nonzero vector $u$ and perturbation $e$ satisfying
$\norm{e}_\infty\le\rho\norm{u}_\infty$ with $\rho<1$,
\[
  \left\|
    \frac{u+e}{\norm{u+e}_\infty}
    -
    \frac{u}{\norm{u}_\infty}
  \right\|_\infty
  \le
  \frac{2\rho}{1-\rho}.
\]
Applying this inequality with $u=\kappa_tq^\star_t$ and
$u=\kappa_tp^\star_t$ proves the claim.
\end{proof}

The stability bounds are the only assumptions specific to the online
recurrence. They say that replacing the current factor maps by the maps used
when each gradient arrived, and replacing the ideal inverse closures by the
online damped inverse closures, changes the weighted quadratic fits only by a
small relative amount.

\paragraph{Relation to KL-Shampoo.}
KL-Shampoo~\cite{klshampoo} also uses the log-determinant divergence to
fit Kronecker factors. The fit here differs in three ways.
\begin{itemize}[leftmargin=1.5em,itemsep=1pt,topsep=2pt]
  \item \emph{No asymptotic distribution limit.} KL-Shampoo fits the second
  moment of a stationary gradient distribution. Here $\EMA_{s\le t}$ is a
  finite weighted sum over realized gradients. The online statement is a
  deterministic perturbation statement for that weighted sum, not an
  infinite-sample limit.
  \item \emph{The fitted moments are factor-gradient moments.} KL-Shampoo
  fits Kronecker factors to the second moment of the weight gradient $G$.
  Here $G$ is never formed. Each diagonal factor is fit from the second
  moment of its corresponding LoRA factor gradient, and the two fits are
  coupled through $C_A$ and $C_B$.
  \item \emph{Fixed-moment statement.} KL-Shampoo
  justifies its moving average as a stochastic proximal gradient step on the
  divergence, correct in the limit. Here \Cref{lem:consistency-p2} is an
  algebraic statement for fixed moments. Under the deterministic moment model,
  using an inexact frozen closure only rescales the fitted vector, which the
  normalization removes. The online EMA is covered by the weighted
  factor-and-closure stability statement in \Cref{prop:online-stationary-p2}.
\end{itemize}

\paragraph{Recovering the update of \Cref{alg:ours}.}
Substituting the fitted model $\kappa_t(Q_t\otimes P_t)$ for
$\overline\Sigma_t$ and the look-ahead momentum approximation for
$\overline G_t$ yields the preconditioned spectral
LMO~\eqref{eq:polorafull} (\Cref{lem:pseudogradient-spectral}). The
direction and magnitude steps of \Cref{sec:sample-loss-lmo}, with the
factor momenta standing in for the two dense momentum projections, give
the update~\eqref{eq:polora-update}. \Cref{alg:ours} updates the
curvature factors after the direction step, so step $t$ uses the estimate
from the previous step. This avoids forming and inverting $C_A$ and
$C_B$ twice per optimizer step.

\subsection{Normalization Placement}
\label{sec:scale-placement-p2}

\methodname{} stores raw accumulated vectors $q_{t+1},p_{t+1}$ and normalizes
them only when forming the effective metrics $D_Q(q_{t+1})$ and
$D_P(p_{t+1})$. It does not normalize each fitted observation
$\hat q_t,\hat p_t$ before adding it to the EMA.

\paragraph{Raw scale does not change the implemented closures.}
For any $a,b>0$,
\[
  D_Q(aq)=D_Q(q),
  \qquad
  D_P(bp)=D_P(p).
\]
The rank-$r$ closures $C_A,C_B$ and the relative-damped inverse closures
$\widetilde C_A^{-1},\widetilde C_B^{-1}$ are therefore unchanged by
positive rescaling of the raw EMA vectors. This is the implementation-level
scale invariance. The magnitude $\kappa_t$ in the diagonal Kronecker
model~\eqref{e-diag-kron-second-mom-model-p2} is removed before the
direction step sees the curvature factors.

\paragraph{Normalizing each observation would change the estimator.}
The recurrence in~\eqref{eq:ours-klcoupling-p2} estimates the weighted
arithmetic mean of the fitted observations,
\[
  q^\circ_{t+1}=\sum_{s\le t}w_{t,s}\hat q_s,
  \qquad
  p^\circ_{t+1}=\sum_{s\le t}w_{t,s}\hat p_s,
\]
with the raw vectors differing by the initialization terms in the unrolled
recurrence. Normalizing before the EMA would replace this by
\[
  \sum_{s\le t}w_{t,s}\frac{\hat q_s}{\norm{\hat q_s}_\infty},
  \qquad
  \sum_{s\le t}w_{t,s}\frac{\hat p_s}{\norm{\hat p_s}_\infty},
\]
which gives a small-magnitude and a large-magnitude fitted observation the
same maximum coordinate before averaging. The implemented order averages the
raw observations first and applies the projective normalization only when the
next effective metric is formed. The same average-then-normalize order appears
in normalized SGD~\cite{cutkosky_nigt}, where momentum precedes
normalization, and in Muon~\cite{denoise_ortho}, where momentum precedes
orthogonalization.
